\documentclass[10pt,a4paper]{article}
\usepackage[a4paper,margin=2.1cm,top=2.4cm,bottom=2.3cm]{geometry}
\usepackage[T1]{fontenc}
\usepackage[utf8]{inputenc}
\usepackage{helvet}
\renewcommand{\familydefault}{\sfdefault}
\usepackage{amsmath,amssymb}
\usepackage{xcolor}
\usepackage{graphicx}
\usepackage{booktabs,tabularx,longtable}
\usepackage{enumitem}
\usepackage{titlesec}
\usepackage{fancyhdr}
\usepackage{tcolorbox}\tcbuselibrary{breakable}
\usepackage[colorlinks=true,linkcolor=cernblue,urlcolor=cernblue]{hyperref}

\definecolor{cernblue}{HTML}{0053A1}
\definecolor{shorange}{HTML}{D55E00}
\definecolor{okgreen}{HTML}{1B7F3B}
\definecolor{greytxt}{HTML}{555555}
\definecolor{lightblue}{HTML}{EAF1F9}

\setlist{itemsep=1pt,topsep=2pt,leftmargin=1.4em}
\titleformat{\section}{\Large\bfseries\color{cernblue}}{\thesection}{0.6em}{}
\titleformat{\subsection}{\large\bfseries\color{cernblue}}{\thesubsection}{0.6em}{}
\titlespacing*{\section}{0pt}{14pt}{6pt}

\pagestyle{fancy}
\fancyhf{}
\renewcommand{\headrulewidth}{0pt}
\fancyfoot[L]{\color{greytxt}\footnotesize SHiP veto optimisation --- issues and possible solutions}
\fancyfoot[R]{\color{greytxt}\footnotesize \thepage}

% status tags
\newcommand{\st}[1]{\textbf{#1}}
\newcommand{\stOPEN}{\textcolor{shorange}{\st{OPEN}}}
\newcommand{\stDONE}{\textcolor{okgreen}{\st{DONE}}}
\newcommand{\stOK}{\textcolor{okgreen}{\st{OK}}}
\newcommand{\stDEF}{\textcolor{greytxt}{\st{DEFERRED}}}
\newcommand{\stPART}{\textcolor{cernblue}{\st{PARTLY DONE}}}
\newcommand{\stSOLVED}{\textcolor{okgreen}{\st{SOLVED}}}

% one issue: #1 id, #2 title, #3 priority, #4 status, #5 source
\newenvironment{issue}[5]{%
  \begin{tcolorbox}[breakable,colback=white,colframe=cernblue,boxrule=0.6pt,arc=1.5mm,
    left=6pt,right=6pt,top=4pt,bottom=4pt,
    title={\textbf{#1}\quad #2\hfill{\normalfont\small priority #3\quad\colorbox{white}{#4}}},
    coltitle=white,colbacktitle=cernblue,fonttitle=\normalsize]
  {\color{greytxt}\footnotesize Source: #5}\par\smallskip}{\end{tcolorbox}\vspace{2pt}}
\newcommand{\prob}{\par\noindent\textbf{\color{cernblue}Issue.}\ }
\newcommand{\sol}{\par\smallskip\noindent\textbf{\color{shorange}Possible solution.}\ }

\begin{document}
\sloppy

% ---------------- title band ----------------
\thispagestyle{fancy}
\begin{tcolorbox}[colback=cernblue,colframe=cernblue,arc=0mm,left=10pt,right=10pt,top=8pt,bottom=8pt]
\color{white}
\begin{minipage}[c]{0.74\linewidth}
{\LARGE\bfseries Veto system optimisation}\\[3pt]
{\Large Issues encountered and possible solutions}\\[6pt]
{\normalsize \underline{G.~Frisella}$^{1,2}$ \quad 4 November 2026}
\end{minipage}\hfill
\begin{minipage}[c]{0.24\linewidth}\raggedleft
\includegraphics[height=1.25cm]{ship-logo-white.png}\\[4pt]
\includegraphics[height=0.75cm]{uzh-logo-white.png}
\end{minipage}
\end{tcolorbox}
{\footnotesize\color{greytxt}$^1$ University of Zurich, $^2$ European Organization for Nuclear Research (CERN)}

\medskip
\noindent This note collects the problems met so far in the fast veto simulation and optimisation
(\texttt{Study\_on\_Veto\_System/Optimization\_Veto\_System/veto}) and, for each, a possible way forward.
The primary source is the open-problem tracker \texttt{ISSUES.txt} (items A1--C9, kept with their IDs, plus A6--A7 added on 4 Nov 2026);
further issues were collected from \texttt{README.md}, \texttt{CHANGELOG.md},
\texttt{docs/SUMMARY\_2026-10-20.md}, \texttt{docs/shield\_iron\_dis.md}, \texttt{docs/fast\_veto\_sim.md},
\texttt{docs/fastveto\_physics.md}, \texttt{docs/geant4\_neutrino\_notes.md} and the FairShip
full-simulation cross-check (items D1--D11). Rates are per spill of $5\times10^{13}$\,PoT
(= kHz at $5\times10^{13}$\,PoT/s); background yields are per 15 years ($6\times10^{20}$\,PoT).

\smallskip
\noindent\textbf{Priority:} A = must be fixed before optimising, B = changes the results, C = second order.
\textbf{Status:} \stOPEN, \stPART, \stDEF, \stSOLVED{} / \stDONE{} / \stOK{} (kept as a record of what was decided).

% ---------------- summary table ----------------
\section*{Overview}
\renewcommand{\arraystretch}{1.12}
\small
\begin{longtable}{@{}l p{0.60\linewidth} c l@{}}
\toprule
ID & Issue & Prio & Status\\\midrule
\endhead
A1 & Pair (combinatorial) background figure of merit saturated & A & \stOPEN\\
A2 & Bounced muons only 74--77\,\% vetoed & A & \stOPEN\\
A3 & Rate limits reported, not enforced & A & \stSOLVED\\
A4 & No objective; two inconsistent cost models & A & \stOPEN\\
A5 & Spill structure and PoT normalisation & A & \stOPEN\\
A6 & Spectrometer field map not used: straight lines from SST1 to the timing detector & A & \stOPEN\\
A7 & No predefined module dimensions; mixed-technology planes & A & \stOPEN\\
B1 & Multiple scattering after the shield & B & \stPART\\
B2 & EM debris (e$^\pm$, $\gamma$) from the shield end not simulated & B & \stOPEN\\
B3 & Neutrino and muon DIS & B & \stPART\\
B4 & Signal model not physical & B & \stOPEN\\
C1 & Detector response simplifications & C & \stOPEN\\
C2 & Dead material and dead regions & C & \stOPEN\\
C3 & No validation against full simulation & C & \stPART\\
C4 & Energy loss of bounced muons in the cavern & C & \stOK\\
C5 & Weights and statistics of rare classes & C & \stOPEN\\
C6 & Detector technology numbers & C & \stDONE\\
C7 & V$^0$ from DIS outside the vessel decaying inside & C & \stOPEN\\
C8 & DIS rates with the new selection & C & \stOPEN\\
C9 & Unreliable bounce flag in the shah\_study categories & C & \stOPEN\\\midrule
D1 & Objective exploitable: signal term saturates at 10\,\% loss & A & \stOPEN\\
D2 & Neutrino DIS in helium: irreducible for an upstream veto & B & \stOPEN\\
D3 & MS7 shield DIS $\to K_L$: chain vs FairShip disagreement & B & \stOPEN\\
D4 & MS7 shield DIS chain: simplifications & C & \stOPEN\\
D5 & Muon source: enriched input gives 1.55$\times$ SST1 rate & B & \stOPEN\\
D6 & Interaction-library approximations & C & \stOPEN\\
D7 & Interactions upstream of 30.6\,m (cavern concrete) not simulated & B & \stOPEN\\
D8 & Results computed with the old PVC vessel wall & B & \stOPEN\\
D9 & DIS$\times\mu$ pairs not redone with the V$^0$-stable library & C & \stOPEN\\
D10 & Geant4 built-in neutrino physics unusable & --- & \stDONE\\
D11 & Bugs found and fixed (record) & --- & \stDONE\\
\bottomrule
\end{longtable}
\normalsize

% =================================================================
\section*{A. Blockers: fix before optimising}

\begin{issue}{A1}{Background figure of merit is saturated}{A}{\stOPEN}{ISSUES.txt A1}
\prob The combinatorial pair background is 2--3 surviving pairs out of 196\,736 candidates
($\sim10^{-5}$ events in 15 years). Every reasonable design gives $\approx0$, so the differences between designs
are Poisson noise and cannot drive an optimiser.
\sol
\begin{itemize}
\item Score designs on the single-muon veto inefficiency per muon type, which has large statistics.
\item Build the pair background by factorisation, $N_{\rm pairs}\approx R_{\rm pairs}\,\varepsilon_1\,\varepsilon_2$, and check
      that it holds on the pairs that do survive.
\item Quote a 90\,\% CL upper limit when there are zero survivors instead of a zero.
\end{itemize}
\end{issue}

\begin{issue}{A2}{Bounced muons: the veto barely sees them}{A}{\stOPEN}{ISSUES.txt A2}
\prob Only 74--77\,\% of the bounced SST1 muons are vetoed. The search window contains the true crossing for only
30\,\% of them: a straight line extrapolated back from the spectrometer cannot follow a bounce off the cavern wall
or floor. This term will probably decide the fin layout. In addition, \texttt{rock\_bounced} and \texttt{veto\_only}
come from a single simulation run (over-fitting risk).
\sol
\begin{itemize}
\item Study this class on its own: where it crosses the planes, before / after the bounce point
      (\texttt{muon\_types.bounce\_points}).
\item Add planes downstream of the bounce point, or a search window aimed at the wall rather than along the
      back-extrapolated line.
\item Produce more independent runs for the bounced classes before optimising on them (see C5).
\end{itemize}
\end{issue}

\begin{issue}{A3}{Rate limits are reported, not enforced}{A}{\stSOLVED}{ISSUES.txt A3; solved 4 Nov 2026 (objective.py, reco.py, bin/make\_peak\_maps.py)}
\prob The 4\,cm tiles on the $z=30.6$\,m fins reach 10\,MHz per channel (limit 1\,MHz; consistent with SR-370's
``up to 10\,MHz in the flames''). The straw fin at 30.6\,m is above 0.5\,MHz/cm$^2$. A penalty $1+\log_{10}(\text{peak}/\text{limit})$
existed, but grew only logarithmically above the limit, and the fin rate maps have $20\times20$\,cm$^2$ bins, which
average the peak down for smaller cells.
\sol \textbf{Done.}
\begin{itemize}
\item Steep penalty in the objective: each limit exceeded costs $w_{\rm rate}\,[1+k\,(\text{peak}/\text{limit}-1)]$,
      $w_{\rm rate}=1000$ (more than all other terms together), $k=10$: a design $2\times$ over its limit costs 11\,000 points.
\item Fine peak maps: $2\times2$\,cm$^2$ maps from all 506\,M muons of the full dataset (UBT at 30.6\,m, fins at
      $z=30.6$--$40.6$\,m, $u<2$\,m); peak over a sliding window of the channel size (pads) or 4\,cm (strips), only windows
      with $\ge10$ effective muons; same normalisation as the coarse maps (0.1\,\%).
\item Effect at 30.6\,m: peak $\times1.39$ (right fin $880\pm5$ vs 634\,kHz/cm$^2$); 4\,cm tiles 14.1\,MHz, 2\,cm tiles
      3.6\,MHz per channel. Penalties: D6\_straw 1\,103 $\to$ 12\,435, D6\_tile 7\,578 $\to$ 253\,426; E3 / E4 stay feasible.
\item Left over (not A3): fibre and pixel have no rate limit yet; fins beyond 40.6\,m use the coarse peak.
\end{itemize}
\end{issue}

\begin{issue}{A4}{No objective, two inconsistent cost models}{A}{\stOPEN}{ISSUES.txt A4}
\prob \texttt{optimize\_sbt.py} is a stub. The cost per layer-m$^2$ (P-369 table 11) and the relative cost per
channel (straw = 1) disagree: fibres are expensive per area, straws cheapest per channel. The pixel cost is
\texttt{nan} in the area model and would break an optimiser. Free parameters and fixed limits are not
written down, and Monte Carlo noise between designs is not controlled.
\sol
\begin{itemize}
\item Fix the objective, e.g.\ minimise cost (as \% of the budget) subject to background $<X$ and signal loss $<Y$
      (see also D1).
\item Keep one cost model (area + channels combined), with a finite pixel cost.
\item Write the geometric limits as constraints: cavern wall $x=-4.56$\,m, 1\,m straw length, nothing below
      $y=-3.2$\,m.
\item Common random numbers: fixed seed and the same events for every design.
\end{itemize}
\end{issue}

\begin{issue}{A5}{Spill and PoT normalisation}{A}{\stOPEN}{ISSUES.txt A5}
\prob Accidental hits use flat times inside a 1\,s spill (\texttt{--spill}); \texttt{config.py} has
$5\times10^{13}$\,PoT per spill and $6\times10^{20}$\,PoT in 15 years. The real spill has a 50\,Hz ripple and a
duty factor $<1$; accidentals scale with the instantaneous rate, and so does the straw signal loss
(5.3\,\% with a $\pm25$\,ns window).
\sol
\begin{itemize}
\item Check that all rates follow the $5\times10^{13}$\,PoT/s convention.
\item Sample hit times from the measured spill profile (ripple + duty factor) and quote the signal loss at the
      peak instantaneous rate.
\end{itemize}
\end{issue}

\begin{issue}{A6}{Spectrometer field map not used}{A}{\stOPEN}{code (muon\_types.py, reco.py, geometry.py), added 4 Nov 2026}
\prob Downstream of SST1 ($z=82$\,m) every track --- muons, DIS products, V$^0$ and signal daughters --- is
extrapolated in a \emph{straight line} to the timing detector ($z=92.71$\,m), through the spectrometer magnet
($z=84.36$--$89.56$\,m). The field map \texttt{MainSpectrometerField.h5} ($z=77$--$97$\,m) is only used for
drawing. The timing-detector acceptance, the $\le2$-tracks cut and every rate after the timing cut therefore
ignore the bending: soft tracks swept out of (or into) the $2\times3$\,m timing plane are counted wrongly, and
opposite-charge pairs do not separate. The $p>1$\,GeV cut uses the true momentum.
\sol
\begin{itemize}
\item Propagate every charged product from its state at $z=82$\,m to the timing plane through the field map
      (vectorised RK4, charge from the PDG code); neutrals stay straight.
\item Use the propagated position for the timing acceptance and the track-multiplicity cuts.
\item Optionally measure $p$ from the bending (stations before / after the magnet, resolution model) and cut on it.
\item Redo the SST1 / TD / $p>1$\,GeV rates (30.6 / 15.3 / 8.5\,kHz) and compare with FairShip
      (48.5 / 14.6 / 7.1\,kHz, field included).
\end{itemize}
\end{issue}

\begin{issue}{A7}{Detector planes have no predefined module dimensions}{A}{\stOPEN}{design.py, technology.py; module sizes from G.~Frisella, 4 Nov 2026}
\prob Plane sizes are free and continuous (fin length, $y$ range, UBT extent); channels and cost are computed for
exactly that area. Only the straw (1\,m) and fibre (3\,m) lengths are quantised. A real plane is assembled from fixed
modules, so the optimiser can propose sizes that cannot be built, and module borders are ignored.
Modules to use (each carries all layers of its technology):

\smallskip
\begin{tabular}{@{}lll@{}}
\toprule
technology & module & content\\\midrule
straw & $1\times1$\,m$^2$ & 6 layers (3 views $0,\pm15^\circ$), 5\,mm pitch\\
tile2 & $50\times50$\,cm$^2$ & $25\times25=625$ tiles of $2\times2$\,cm$^2$\\
tile & $48\times48$\,cm$^2$ & $12\times12=144$ tiles of $4\times4$\,cm$^2$\\
fibre & $0.5\times2.5$\,m$^2$ & two crossed SCSF-78MJ mats (x, y), 250\,$\mu$m\\
pixel & $\sim4\times4$\,cm$^2$ & quad of $2\times2$ Timepix4-like chips\\
\bottomrule
\end{tabular}
\smallskip

\noindent Modules of every technology are assembled the same way: side by side (an $n_u\times n_y$ grid, e.g.\ 4 straw
modules $=2\times2$\,m$^2$) and stacked in $z$ at the same $(x,y)$ (redundant stations). Technologies can be
\emph{mixed} in one plane: e.g.\ 1\,m in $x$ = two $50\times50$\,cm$^2$ tile2 modules, or one tile2 module + one
0.5\,m fibre mat, so fine timing / granularity goes only where the rate needs it. Each module has a dead frame
(frame, SiPMs, readout). The module sizes are not all on one grid (tile 48 vs 50\,cm, fibre 2.5\,m vs straw 1\,m in
$y$): a common grid (e.g.\ 50\,cm, tile modules padded) or mismatched borders as dead area must be decided.
\sol
\begin{itemize}
\item \texttt{technology.py}: module size, cells per module, dead-frame width (to be fixed with the detector
      colleagues), channels and cost per module.
\item \texttt{design.py}: a plane = a grid of module slots, each slot = (technology, number of stacked layers);
      these integers / choices become the optimiser's free parameters.
\item \texttt{reco.py}: zero efficiency in the dead frames and gaps; frame material in the MS / DIS budget (C2).
\item Rate limits (A3) per module channel; redo the existing designs on the module grid.
\end{itemize}
\end{issue}

% =================================================================
\section*{B. Physics missing: changes the results}

\begin{issue}{B1}{Multiple scattering after the shield}{B}{\stPART}{ISSUES.txt B1, docs/fastveto\_physics.md, CHANGELOG 2026-11-04}
\prob The stored muons have no multiple scattering after the shield: the states at $z=30.6$\,m and $z=82$\,m are
joined by straight lines. The material budget is now defined (air, 350\,$\mu$m Mylar entrance / exit walls, 50\,m
helium, the veto planes): $x/X_0=0.0182$ without the UBT, 0.019 with a straw UBT, 0.043 with a tile UBT, i.e.\
$\theta_0\approx1.6$--$2.5$\,mrad$/p$[GeV]. Relevant for the soft muons (1--few GeV) that dominate the fin rates.
A Highland + tails model (\texttt{scattering.py}, \texttt{fastveto.msc}) is implemented and validated against Geant4
on the vessel path; still open:
\begin{itemize}
\item the dense planes (tile, 2.4\,\% $X_0$) are not checked against Geant4;
\item the material of the vessel \emph{side} walls (muons entering through the SBT) is not defined;
\item for bounced muons MS must start only after the bounce point.
\end{itemize}
\sol
\begin{itemize}
\item Run \texttt{bin/msc\_validate.py} with a tile plane in the path.
\item Define the side-wall material and add it to the per-segment walk (A$\to$K$\to$B).
\item Check that \texttt{RecoConfig.ms\_slope} / \texttt{r\_ms} (window growth with $1/p$) still cover the true
      scatter; otherwise the window loses efficiency or gains accidentals.
\end{itemize}
\end{issue}

\begin{issue}{B2}{EM debris from the end of the shield}{B}{\stOPEN}{ISSUES.txt B2}
\prob Only muons are simulated: no $\delta$ rays, bremsstrahlung or pair-production debris leaking from the last
iron of the shield. These add hits (occupancy, rate limits) and accidental vetoes (signal loss). The FairShip
full-physics cross-check shows how large this can be: all charged particles at SST1 are $\sim3.6$\,MHz vs 48.5\,kHz
of muons, dominated by e$^\pm$.
\sol
\begin{itemize}
\item From the full dataset (506\,M rows at $z=30.6$\,m), compute per muon the probability to produce an
      escaping e$^\pm$ / $\gamma$ in the last 20\,cm of iron ($z=29.69$--$29.89$\,m, 11.4\,$X_0$, equilibrium layer):
      $\delta$ rays ($dN/dT\simeq0.56$\,MeV\,cm$^{-1}/T^2$), bremsstrahlung, direct pair production.
\item Turn it into e/$\gamma$ rate maps on the UBT and fins, apply a per-technology detection probability
      (photons rarely convert in thin straws, more in 1\,cm tiles), add to the accidentals.
\item Cross-check with Geant4 (muon beam on 20\,cm iron) and with the FairShip EM sample.
\end{itemize}
\end{issue}

\begin{issue}{B3}{Neutrino and muon DIS}{B}{\stPART}{ISSUES.txt B3, README, docs/fast\_veto\_sim.md}
\prob Originally deferred. Since then forced muon DIS (Geant4 + Pythia 6 library) and neutrino interactions
(GENIE 3.06.02) are modelled in helium, air, vessel walls, every plane and the last 50\,cm of MS7. Still missing:
DIS in the cavern (see D7) and the material term in the optimiser --- without it the optimiser favours thick, fast
planes for free (tiles 2.4\,\% $X_0$ vs straws 0.12\,\%).
\sol
\begin{itemize}
\item Keep the per-plane interaction source in the objective (already in \texttt{objective.score}) so plane
      material is penalised through its own DIS.
\item Update the status in \texttt{ISSUES.txt} (still says DEFERRED).
\end{itemize}
\end{issue}

\begin{issue}{B4}{Signal model}{B}{\stOPEN}{ISSUES.txt B4}
\prob The signal is a neutral particle decaying uniformly in the vessel, two daughters uniform over the SST1
acceptance, $p$ uniform 1--50\,GeV. Signal is lost only through accidental hits. Opening angles, momenta and the
vertex distribution are not realistic.
\sol
\begin{itemize}
\item Use HNL and dark-photon models from the literature: decay probability
      $P=e^{-L_1/\lambda}-e^{-L_2/\lambda}$, $\lambda=\beta\gamma c\tau$; production from charm/beauty at 400\,GeV;
      channels $\mu\pi$, $e\pi$, $\mu\mu\nu$; 2/3-body kinematics. Take them from the FairShip generator or
      [1--5], not invented.
\item A truth-level path exists: the collaboration LLP generator EventCalc (per-channel cut efficiencies already
      studied in \texttt{EventCalc\_signal\_rejection}); feed its decays through the veto.
\end{itemize}
{\footnotesize\color{greytxt}[1] SHiP, JHEP 04 (2019) 077; [2] Bondarenko et al., JHEP 11 (2018) 032;
[3] SHiP, EPJC 81 (2021) 451; [4] Alekhin et al., Rep.\ Prog.\ Phys.\ 79 (2016) 124201; [5] CERN-SPSC-2023-033.}
\end{issue}

% =================================================================
\section*{C. Detector and bookkeeping: second order}

\begin{issue}{C1}{Detector response simplifications}{C}{\stOPEN}{ISSUES.txt C1}
\prob $\sigma_x$ is not used (the 5\,cm search window dominates); no ghost hits from combining views when several
muons share a time window; straw time = uniform smearing over the 50\,ns drift window; no cluster size, no tile
cross-talk; plane efficiencies independent.
\sol Add ghost-hit combinatorics per time window (they matter at the fin rates); with a tile $t_0$ use the
drift radius for the straw time and position; correlated inefficiencies (module-level) as a systematic.
\end{issue}

\begin{issue}{C2}{Dead material and dead regions}{C}{\stOPEN}{ISSUES.txt C2}
\prob No frames, manifolds, overlaps between 1\,m straw rows, module gaps, SiPMs or cables. For the UBT this
matters most: the tensioned-straw frame sits in the acceptance (SR-370 open question).
\sol Add the UBT frame and module gaps as passive material (MS + DIS target) and as dead regions (zero efficiency);
measure the inefficiency of the gaps with the single-muon figure of merit (A1).
\end{issue}

\begin{issue}{C3}{Validation against full simulation}{C}{\stPART}{ISSUES.txt C3, FairShip EM SST1 cross-check}
\prob The fast simulation was never run end-to-end against FairShip for one design. A first full-physics check
(FairShip \texttt{TRY\_2026\_EM\_SST1\_origin}, 411.6\,M muons) gives muons at SST1 / TD / $p>1$\,GeV of
48.5 / 14.6 / 7.1\,kHz vs 30.6 / 15.3 / 8.5\,kHz in the fast chain; the excess after the first cut is
$\pi$ decays in the vessel. Shield-made hadrons disagree strongly (see D3).
\sol Run one baseline design through FairShip and compare rates per plane and the veto inefficiency per
muon type before trusting an optimum.
\end{issue}

\begin{issue}{C4}{Energy loss of bounced muons in the cavern}{C}{\stOK}{ISSUES.txt C4}
\prob Whether the energy loss in the cavern rock is modelled for the bounced muons.
\sol Nothing to do: it is in the stored muons from the transport.
\end{issue}

\begin{issue}{C5}{Weights and statistics}{C}{\stOPEN}{ISSUES.txt C5, CHANGELOG 2026-10-30}
\prob The repeated-\texttt{evt\_n} weight correction is not applied in the optimisation dataset;
\texttt{veto\_only} and \texttt{rock\_bounced} come from one run. For muon DIS in helium one input muon gives
$6.8\times10^4$ of $10^6$ draws: $n_{\rm eff}\approx60$, $\pm23$\,\% on the rate. Flattening the weights to 1 was
rejected: the $w=768.75$ minimum-bias class is 40.5\,\% of the true shield DIS rate.
\sol Decide on the \texttt{evt\_n} correction; one forced interaction per crossing / per track is already used
(errors halved for $\nu$); produce more independent seeds for the rare classes, and keep bootstrapping over
\texttt{evt\_n} for all errors.
\end{issue}

\begin{issue}{C6}{Detector technology numbers}{C}{\stDONE}{ISSUES.txt C6 (commit f6d13ec)}
\prob Straw / fibre / pixel / tile parameters needed the colleagues' input.
\sol Updated in \texttt{technology.py}. Still to confirm: fibre timing (0.5\,ns kept, may reach $\sim100$\,ps),
per-SiPM efficiency, tile channel count for mixed tile sizes.
\end{issue}

\begin{issue}{C7}{V$^0$ from DIS outside the vessel decaying inside}{C}{\stOPEN}{ISSUES.txt C7}
\prob A $K_S$ / $\Lambda$ made by muon DIS in cavern air, a fin or the UBT can decay inside the vessel with no
upstream hit of its own. Air: 0 of 45 air pairs after DOCA have a vertex in the vessel, but one air draw
corresponds to $7\times10^6$ interactions per 15\,y (not conclusive). Fins / UBT are not in the interaction caches;
$K_L$ kept stable.
\sol Draws restricted to air and fins close to the vessel entrance (\texttt{dis\_cutflow.py}); add the fins / UBT
to the interaction caches; decay $K_L$ in flight (as already done for MS7, D3).
\end{issue}

\begin{issue}{C8}{DIS rates with the new selection}{C}{\stOPEN}{ISSUES.txt C8}
\prob Without the measured IP $<10$\,cm and $\le2$ tracks in SST1, before the veto per 15\,y:
$\mu$ He $(6.0\pm1.4)\times10^4$, $\mu$ wall $7\,400\pm850$ (old PVC), $\nu$ He $467\pm2$
(old selection: 1\,800 and 0.82).
\sol Decide whether to bring back a pointing / track-multiplicity cut; the daughter-asymmetry cuts
($|A_p|,|A_{IP}|<0.9$) keep 48\,\% of $\mu$ He DIS and 99\,\% of the wall DIS, so they alone are not enough.
Redo the wall numbers with Mylar (D8).
\end{issue}

\begin{issue}{C9}{shah\_study categories: bounce flag not reliable}{C}{\stOPEN}{ISSUES.txt C9, CHANGELOG 2026-11-04}
\prob In \texttt{cuda\_muons\_ship.categorize\_muons} every \texttt{bounced\_back} muon checked (seed 0) is a perfect
straight line from 30.6\,m to 82\,m, and \texttt{straight\_past\_vessel} excludes straight muons through the vessel
and contains 0.6\,\% bounced ones. With the category the straight \& selected UBT rate was 6.24\,kHz instead of
8.13\,kHz.
\sol Use the veto definition (offset from the 30.6\,m line $<$ \texttt{BOUNCE\_CUT} = 0.5\,m,
\texttt{muon\_types.py}); report the bug to the MuonsAndMatter authors.
\end{issue}

% =================================================================
\section*{D. Further issues from the other documents}

\begin{issue}{D1}{Objective exploitable: signal term saturates}{A}{\stOPEN}{docs/SUMMARY\_2026-10-20.md \S5--6}
\prob $\text{reward}=-[100(1-e^{-N})+10\min(\text{sig.\ lost}/0.10,1)+\text{cost}]$. The signal term saturates at
10\,\% loss, so losing 10\,\% and 74\,\% cost the same. Once neutrinos dominate $N$, vetoing everything pays:
evolution run 3 reached its best reward with a design losing 74\,\% of the signal (50\,mrad cone, $\sim15$ planes)
to bring $\nu$ from 0.6 to 0.06.
\sol
\begin{itemize}
\item A sensitivity figure of merit, e.g.\ maximise $\varepsilon_{\rm sig}/N_{90\%\,\rm UL}$.
\item Leave neutrinos out of $N$ in the optimiser (or count them before the veto), since an upstream veto
      cannot remove them (D2).
\item A linear, uncapped signal term with a hard penalty above $\sim5$\,\% loss.
\end{itemize}
\end{issue}

\begin{issue}{D2}{Neutrino DIS in helium: irreducible for an upstream veto}{B}{\stOPEN}{docs/SUMMARY\_2026-10-20.md \S4}
\prob About 0.6 $\nu$ DIS in the helium per 15\,y survive all designs: the vertex is in the gas, nothing charged
comes in, and the third that is vetoed is vetoed by random hits (like signal). Tagging the event's own extra
products (wide-angle hadrons hitting fins downstream of the vertex, SBT-style) is not modelled:
\texttt{interactions.charged\_tracks} keeps only forward tracks and the veto searches only upstream.
\sol Keep all charged / neutral final-state particles, digitise hits downstream of the vertex on the fins, and add
a ``vertex-activity'' veto; evaluate its signal cost with the realistic signal model (B4).
\end{issue}

\begin{issue}{D3}{MS7 shield DIS $\to K_L$: chain disagrees with FairShip}{B}{\stOPEN}{docs/shield\_iron\_dis.md, FairShip EM SST1 check}
\prob The forced chain in the last 50\,cm of MS7 gives $K_L$ decays in the vessel of $\sim4.9$\,kHz and
$(6.73\pm0.59)\times10^7$ per 15\,y after DOCA, fiducial and IP $<2.5$\,m. Concerns:
\begin{itemize}
\item unforced Geant4 gives only $0.52\pm0.15$ of the chain's $K_L$; the excess comes from the Pythia 6
      component of \texttt{mu\_library\_iron} (9.5\,\% of interactions, 90\,\% of the $K_L$);
\item FairShip full physics sees \emph{zero} $K_L$/$K_S$ from the shield reaching the SST, and shield-made
      $\pi$ are $\sim17\times$ below the chain --- not explained by the Pythia excess;
\item after the full selection the $\mu$ number rests on 3--7 MC events ($\pm40$--60\,\%);
\item algorithm A cannot veto it (decay tracks have no upstream hits); asymmetry cuts keep 100\,\%.
\end{itemize}
\sol Find the source of the FairShip discrepancy: MS7 flux and field TRY6 vs TRY2026, the muon-nuclear model
(Kokoulin in FairShip vs Pythia 6 high-$\nu$ DIS here), the flux footprint already agrees. Use $K_L$-enriched
draws for statistics. Model a veto on any in-time hit (scattered muon in the UBT / fins) and a pointing cut.
\end{issue}

\begin{issue}{D4}{MS7 shield DIS chain: simplifications}{C}{\stOPEN}{docs/shield\_iron\_dis.md (Limitations)}
\prob Straight parent lines from 29.20\,m (no bending in 2\,T, $\sim1.4$\,cm at 20\,GeV); muons stopping in the
slab not removed; Pythia part on free nucleons; straight lines and no field after the plane; $\pi/K$ decays in
flight ignored; neutron and $K_L$ interactions in the vessel / air not simulated; neutrino flux $p>10$\,GeV only,
no $\nu_e$.
\sol Propagate the parent with the field map (already in \texttt{shield\_iron.py}); decay $\pi/K$ in flight
and transport neutrons / $K_L$ in helium and air; extend the neutrino flux below 10\,GeV and add $\nu_e$.
\end{issue}

\begin{issue}{D5}{Muon source: enriched input disagrees with the old sample}{B}{\stOPEN}{docs/fast\_veto\_sim.md \S5}
\prob The SST1-selected rate is $1.55\times$ higher with the enriched input ($1.9\times$ above 100\,GeV) at the
same flux above 150\,GeV --- larger than any statistical error. In addition the regeneration differs by
10--20\,\% from the old file, the $>200$\,GeV tail has a $\sim\times2$ model dependence (Pythia version /
diffraction), and 50--150\,GeV still has the old statistics.
\sol Compare angle and momentum shapes of the two inputs above 150\,GeV; extend the enrichment down to
50\,GeV; quote the source model as a systematic on the muon backgrounds.
\end{issue}

\begin{issue}{D6}{Interaction-library approximations}{C}{\stOPEN}{docs/fastveto\_physics.md (Open)}
\prob Final states are not rescaled to the exact muon energy (15\,\% grid); no nuclear effects (Fermi motion,
shadowing, re-interactions) in the Pythia part; no $\tau$ library ($\sim1$\,\% $\nu_\tau$ in the flux).
\sol Rescale final states to the true energy (or finer grid); compare with a Geant4/GENIE nuclear model on
a subset; build a $\nu_\tau$ library.
\end{issue}

\begin{issue}{D7}{Interactions upstream of 30.6\,m not simulated}{B}{\stOPEN}{docs/fast\_veto\_sim.md \S5, FairShip origin breakdown}
\prob Apart from the last 50\,cm of MS7, interactions in the shield, cavern walls and rock are not simulated.
FairShip shows that $\pi$ / p at SST1 come mostly from the cavern concrete and the vessel walls / SBT.
\sol Add the cavern concrete (and the rest of the shield) as forced-DIS targets using the bounced-muon
paths; use the FairShip origin breakdown as the reference.
\end{issue}

\begin{issue}{D8}{Results computed with the old PVC vessel wall}{B}{\stOPEN}{CHANGELOG 2026-11-04}
\prob The vessel wall is now 350\,$\mu$m Mylar ($x/X_0$ 1.53 $\to$ 1.23\,$\times10^{-3}$, +18\,\% nucleons, no Cl).
\texttt{results/interactions}, apply / reco-input files with ``pvc'' rows, scattering-dependent reco and the C8
wall numbers are still PVC-based.
\sol Rerun the wall DIS and the reco with Mylar; rename the target key \texttt{"pvc"} to \texttt{"wall"} once all
saved files are regenerated.
\end{issue}

\begin{issue}{D9}{DIS$\times\mu$ pairs not redone}{C}{\stOPEN}{shield DIS notes (2026-11-02)}
\prob With $K_S$/$\Lambda$ now tracked in the MS7 iron (\texttt{mu\_library\_iron\_v0stable.npz}), the
\texttt{dis\_comb\_*} DIS$\times\mu$ combinatorial pairs were not redone.
\sol Rerun \texttt{dis\_comb\_pairs\_asym.py} on the V$^0$-stable sample.
\end{issue}

\begin{issue}{D10}{Geant4 built-in neutrino physics unusable}{---}{\stDONE}{docs/geant4\_neutrino\_notes.md}
\prob In Geant4 11.2.2 all $\nu$--nucleus models share one set of $x/Q^2$ tables (CC sampled from NC
kinematics, $\langle y\rangle=0.13$--0.34 instead of 0.48), NC is $2.6\times$ too large and drops all hadrons,
no Lund fragmentation, and a \texttt{fWeight} bug gives CC secondaries weight 1.
\sol Done: GENIE 3.06.02 (tune G18\_02a) built and used for all neutrino interactions.
\end{issue}

\begin{issue}{D11}{Bugs found and fixed (record)}{---}{\stDONE}{CHANGELOG, SUMMARY 2026-10-20}
\noindent Bugs that changed results, with the fix:
\begin{itemize}
\item \texttt{reco.own\_hits}: with MS on, tracks born downstream got a UBT hit at their start (9.5\,\% fake
      vetoes for $\nu$ DIS in helium) --- start-upstream condition restored (2026-11-01).
\item \texttt{make\_design\_explorer.py}: muons stopping in the shield (type 7) counted on the UBT
      (5.3$\times10^9$ instead of 1.05$\times10^7$ per spill) --- excluded.
\item \texttt{shield\_dis}: muon final states sampled twice per interaction --- fixed.
\item shah\_study ``straight'' category replaced by the \texttt{BOUNCE\_CUT} definition (C9).
\end{itemize}
\end{issue}

\end{document}
